% SGA 6, Exposición II: fragmento español alineado con la fuente.
% Alcance: sección 3, demostración de la proposición 3.7 b), líneas 4310--4334 inclusive.
% Fuente francesa: sga6_fr_r1/src/fr.tex.
% SHA-256 de la fuente francesa completa:
% 0D2158B9FE3B466309A92165727A0B191422E05FBAA0DACD5806E8322C54BA1A.
\providecommand{\sisi}[2]{#2}

\emph{Demostración de b).} Como el coherizador es de dimensión cohomológica finita según a), admite ([H] I 5.3) un derivado derecho
\[
 RQ:\mathrm D(S)\longrightarrow\mathrm D(\Qcoh(S))
\]
(cuya restricción a $\mathrm D^+(S)$ coincide con el funtor $\mathbb RQ$ definido anteriormente). Para $E\in\Ob\mathrm D(\Qcoh(S))$, existe un isomorfismo funtorial canónico
\begin{equation}\tag{1}
 E\xrightarrow{\sim}\mathbb RQ(E),
\end{equation}
puesto que, según (3.5.1), los $E^i$ son acíclicos para $Q$ y tales que $E^i\xrightarrow{\sim}Q(E^i)$. Por otra parte, para $E\in\Ob\mathrm D(S)$, se tiene un homomorfismo funtorial canónico
\begin{equation}\tag{2}
 RQ(E)\longrightarrow E
\end{equation}
(que proviene, cuando los $E^i$ son $Q$-acíclicos, del homomorfismo de adjunción ordinario entre $\varphi$ y $Q$). Se verifica trivialmente que los homomorfismos (1) y (2) hacen de
\[
 RQ:\mathrm D(S)\longrightarrow\mathrm D(\Qcoh(S))
\]
un adjunto a derecha de
\[
 \varphi:\mathrm D(\Qcoh(S))\longrightarrow\mathrm D(S).
\]
El hecho de que (1) sea un isomorfismo implica que $\varphi$ es plenamente fiel. Finalmente, utilizando la sucesión espectral
\[
 R^pQ(H^q(E))\Longrightarrow R^{*}Q(E),
\]
(convergente porque $Q$ es de dimensión cohomológica finita), se ve que (2) es un isomorfismo para $E\in\Ob\mathrm D(S)_{\qcoh}$, y por tanto que $\varphi$ tiene por imagen esencial $\mathrm D(S)_{\qcoh}$. Esto prueba b) y termina la demostración de (3.7).
